% 原书分章PDF第1—3页：章首与第30.1节。
\ctexset{chapter/aftername=\par}
\chapter{真理的困境：论统计学、因果性、\mbox{物理学与科学哲学}}
\ctexset{chapter/aftername=\quad}
\label{ch:predicament-of-truth}
\chapterauthors{Richard J. C. M. Starmans}
\chapteraffiliations{%
R. J. C. M. Starmans（通讯作者）：荷兰乌得勒支大学信息与计算科学系，Buys Ballot 实验室，Princetonplein 5, De Uithof, 3584 CC Utrecht, The Netherlands；电子邮箱：\href{mailto:starmans@cs.uu.nl}{starmans@cs.uu.nl}。\\
本章 DOI：\url{https://doi.org/10.1007/978-3-319-65304-4_30}。}

\section{统计学与真理的脆弱性}

% S30.1.P01
一方面，本文标题的前半部分或许显得有些怨怼、不祥，仿佛是一连串抱怨的开场，或是为追求真理之路的衰败而悲伤、痛苦地沉思的前奏。不可否认，它暗示着我们所熟悉的文明将进一步衰落，呼应了奥斯瓦尔德·斯宾格勒（Oswald Spengler）关于文明寿命的悲观预言；他的《西方的没落》（Der Untergang des Abendlandes）于第一次世界大战结束的1918年出版。真理概念在思想史中始终至关重要，也构成了人类处境的特征与独特之处。即使以温和或激进的方式信奉普罗泰戈拉（Protagoras）的“人是万物的尺度”原则（homo-mensura-principle），人们也不可能在不宣称、规定、猜想、确立或珍视某种真理观念的情况下存在、生存或行事；这种观念支撑并激励着他们的思想、言语和行动。因此，真理在宗教、形而上学（metaphysics）、认识论（epistemology）、科学、政治和日常生活中无处不在。从历史哲学的角度看，真理概念在认识论中处于\emph{枢纽地位}：它先于、涵摄，或至少是知识、理性、客观性、因果性、证成（justification）、推断等许多概念所\emph{预设}的概念。与此同时，人们又很容易以自身利益、宗教、政治意识形态、自由、利益攸关方的利益、公共卫生、国家安全、气候、人民意志等为名，否认真理、歪曲真理、宣布真理已经过时，或为行方便而修改真理、使之相对化。在政治、新闻业——无论表现为虚假信息、被宣称为不言自明的道理，还是假新闻——以及社交媒体、历史编纂中，都可以看到真理的这种脆弱性；甚至，说来可怕（horribile dictu），在哲学与现代科学中也不例外。

% S30.1.P02
另一方面，资深统计学家或许不会对上述图景感到如此惊讶或震恐，因为自这门学科诞生以来，他们就一直因统计学与真理之间充满问题的关系而受到审判。马克·吐温（Mark Twain）那句臭名昭著的格言“有谎言，有该死的谎言，还有统计”，以及乔治·博克斯（George Box）那句标志性的警语“所有模型都是错的，但有些是有用的”，都曾为各种目的而被频繁引用。Darrell Hoff 的小册子《如何用统计学说谎》（How to Lie with Statistics，1954）在全球售出超过一百万册，甚至催生了一类专门讨论统计谬误、警示与陷阱的新作品。著名的统计史著作，例如 Stigler（1986，1999，2016）、Porter（1986，1995）、Krüger 等（1987）、Hacking（1975，1990）、Gigerenzer（1989）的著作以及许多其他作品，都充分表明：在统计学发展的整个过程中，无论是在科学、哲学还是大众文化中，怀疑与不信任始终普遍存在。

% S30.1.P03
例如，伊恩·哈金（Ian Hacking）在其专著《偶然性的驯服》（The Taming of Chance，1990）中，描绘了19世纪一股持续而强韧的反统计学潮流。持这种态度的重要知识分子来自不同阵营，包括陀思妥耶夫斯基（Fjodor M. Dostoyevsky）、奥古斯特·孔德（August Comte）和弗里德里希·尼采（Friedrich Nietzsche）。他们反对统计学，因为它（被认为）倚重效用、数值抽象、平均值与平庸，而不去追求真理、真正的知识和理解，不去关注本质，也不去维护每一个人的尊严。此后，统计学领域虽然取得了显著进展，而且几乎所有科学都在上个世纪发生了概率转向，这是无可争议的事实，但科学、政治和大众文化中的怀疑仍然持续存在；即使在今天，伦理学与统计学之间的对话也往往艰难而费力（Starmans，2012a,b,c）。

% S30.1.P04
当然，这些简短引语和文献引用，只反映了真理与统计学关系的少数几个方面。这里不讨论统计这一职业在多大程度上应为大众文化中的上述看法负责。我们将把范围限于当前数理统计的研究，并提出三点观察。

% S30.1.P05
第一，从历史哲学的角度看，靶向最小损失估计（targeted minimum loss-based estimation，TMLE）范式的目标之一，显然是重新思考，或者更确切地说，恢复真理与模型之间的联系。它以审慎而简约的方式重新确立统计模型的概念，只允许人们把自己真实、符合实际的知识纳入模型，例如数据经过随机化，代表某一随机变量的独立同分布观测。TMLE 不去假定设定错误的参数统计模型，或限制过强的半参数统计模型，而是用非参数或半参数统计模型来定义统计估计问题；这些模型表示符合实际的知识，即关于实际产生手头数据的实验的知识。模型必须包含真实概率分布（的最优逼近），而靶向的感兴趣参数则定义为这一真实分布的函数。为此，TMLE“将估计这一概念本身——它历来在统计推断中具有经典地位——重新置于任何关于从数据学习的理论或方法的主导位置”（Starmans，2011）。

% S30.1.P06
第二，当前统计数据分析的实践既充满悖论，也充满问题。一方面，教科书错误地让人以为，这是一片基础清晰、技术有充分依据的统一领域；它们叙述的是一种传统的\emph{辉格式历史}（whig history），其中知识不断产生、积累并得到证成。另一方面，我们在现实中看到的却是大量不同、有时甚至互不相容的技术与原则；要运用它们，必须经过细致的\emph{折衷选择}（eclecticism）。这些技术与原则往往已经脱离其历史背景，而且源自有时互不相容或彼此不协调的立场，包括 Karl Pearson、Udny Yule、Ronald Fisher、Jerzy Neyman、Egon Pearson、Jimmy Savage、Bruno de Finetti 和 John Tukey 的立场（Starmans，2011）。

% S30.1.P07
第三，数据科学近期的发展，以及深度学习等各类强大的机器学习技术，很可能进一步恶化这一局面。统计学与算法数据分析之间的协调似乎势在必行；实际上，促成这种协调正是 TMLE 的指导原则之一（van der Laan 和 Rose，2011；van der Laan 和 Starmans，2014）（Starmans，2016a）。本文第三节将进一步讨论这一问题，重点考察当前在数据科学中日益流行的\emph{反统计学立场与反认识论立场}（anti-statistical and anti-epistemic stance）。
